% SEGMENT OLP-0483-S01
% Part: applied-modal-logics
% Chapter: epistemic-logic
% Section: introduction

\documentclass[../../../include/open-logic-section]{subfiles}

\begin{document}

\olfileid{aml}{el}{int}

\olsection{அறிமுகம்}

வாய்ப்புநிலைத் தருக்கம் \emph{வாய்ப்புநிலைக் கூற்றுகளையும்} அவற்றுக்கிடையிலான
பின்விளைவு உறவுகளையும் கையாள்வதுபோலவே, அறிவுநிலைத் தருக்கம்
\emph{அறிவுநிலைக் கூற்றுகளையும்} அவற்றுக்கிடையிலான பின்விளைவு உறவுகளையும்
கையாள்கிறது. வாய்ப்புநிலைச் செயற்குறிகள் கூடுதன்மையையும் கட்டாயத்தன்மையையும்
குறிப்பதாகப் பொருள்கொள்வதற்குப் பதிலாக, அறிவையும் நம்பிக்கையையும் மாதிரிப்பதற்கு
ஓருறுப்பு இணைப்பிகள் அறிவுநிலையாகவோ நம்பிக்கைநிலையாகவோ
பொருள்கொள்ளப்படுகின்றன. எடுத்துக்காட்டாக, பின்வருவன போன்ற கூற்றுகளை நாம்
வெளிப்படுத்த விரும்பலாம்:

\begin{enumerate}
\item கல்கரி ஆல்பர்ட்டாவில் இருப்பதை ரிச்சர்ட் அறிகிறார்.
\item ஒரு நாய் சோபாவில் இருப்பது கூடும் என்று ஆட்ரி நினைக்கிறார்.
\item தனது வகுப்பு செவ்வாய்க்கிழமைகளில் நடைபெறுவதை ஆட்ரி அறிவதை ரிச்சர்ட் அறிகிறார்.
\item ஓராண்டில் 12 மாதங்கள் உள்ளன என்று எல்லோரும் அறிகிறார்கள்.
\end{enumerate}
தற்கால அறிவுநிலைத் தருக்கத்தின் தொடக்கம், 1962 இல் வெளியான யாக்கோ
ஹின்டிக்காவின் \emph{அறிவும் நம்பிக்கையும்} என்ற நூலுக்கு அடிக்கடி
இட்டுச்செல்லப்படுகிறது; தருக்கத்தில் சாத்தியமான உலகப் பொருண்மையியல் அதிகமாகப்
பயன்படத் தொடங்கிய காலத்தில் அந்நூல் எழுதப்பட்டது. உண்மையில், பிற வாய்ப்புநிலைத்
தருக்கங்களின் பொருண்மையியல் கருவிகளில் பெரும்பாலானவற்றையே அறிவுநிலைத்
தருக்கங்களும் பயன்படுத்துகின்றன; ஆனால் அவற்றை வேறுவிதமாகப் பொருள்கொள்கின்றன.
\emph{அணுகுறவு} எதைக் குறிக்கிறது என்று நாம் கொள்வதில்தான் முதன்மையான மாற்றம்
உள்ளது. அறிவுநிலைத் தருக்கங்களில், அவை ஏதோவொரு வகை \emph{அறிவுநிலைக்
கூடுதன்மையை} குறிக்கின்றன. நாம் மாதிரிக்கும் அறிவுநிலைக் கருத்து, அணுகுறவின்
மீது விதிக்க விரும்பும் கட்டுப்பாடுகளை எவ்வாறு பாதிக்கிறது என்பதைக் காண்போம்.
சட்டகப் பொருத்தக் கோட்பாட்டிற்கு அறிவுநிலைப் பொருள்கோள் அளிக்கப்படும்போது என்ன
நிகழ்கிறது என்பதையும் காண்போம். மேலுள்ள எடுத்துக்காட்டுகள் ரிச்சர்ட், ஆட்ரி
என்ற இரண்டு முகவர்களையும், அவர்கள் ஒவ்வொருவரும் அறிந்தவற்றுக்கிடையிலான
உறவையும் குறிப்பிடுவதைக் கவனிக்கவும். நாம் ஆராயவுள்ள அறிவுநிலைத் தருக்கங்கள்,
இத்தகையவற்றை வெளிப்படுத்தக்கூடிய பன்முகவர் தருக்கங்களாக இருக்கும். இதற்கு
மாறாக, ஒருமுகவர் அறிவுநிலைத் தருக்கம் ஓர் ஆள் அறிந்த அல்லது நம்பும்
விடயங்களைப் பற்றி மட்டுமே பேசும்.

\end{document}
